% Copyright (c) 2026 by C. Klukas.
% Licensed under the MIT License. See LICENSE for details.

A mesh is a set of hosts that trust the same signed ledger. Nothing else
binds them: there is no central server, no shared secret, and no
membership list or policy outside the ledger. This section describes the ledger and the identities
in it, the passports that paglets carry, the policy that decides access to
host resources, the system paglets through which that access is given, the
contracts that define them, and the two services with which a paglet sees
the mesh and gets compute time on it. Figure~\ref{fig:mesh} shows the
parts.

\begin{figure}[htbp]
\centering
\begin{tikzpicture}[node distance=9mm and 12mm]
  \node[box] (h1) {host \code{lab-1}\\[1pt]\footnotesize ledger copy, runtime,\\[-2pt]\footnotesize system paglets};
  \node[box, right=of h1] (h2) {host \code{lab-2}\\[1pt]\footnotesize ledger copy, runtime,\\[-2pt]\footnotesize system paglets};
  \node[box, right=of h2] (h3) {host \code{laptop}\\[1pt]\footnotesize ledger copy, runtime,\\[-2pt]\footnotesize system paglets};
  \draw[both] (h1) -- node[event, above] {gossip} (h2);
  \draw[both] (h2) -- node[event, above] {gossip} (h3);
  \node[plain, above=of h1] (admin) {admin\\\footnotesize admin key};
  \draw[arrow] (admin) -- node[event, right] {signed records} (h1);
  \node[plain, above=of h3] (owner) {owner\\\footnotesize owner key};
  \draw[arrow] (owner) -- node[event, left] {module, passport} (h3);
  \node[state, below=of h2] (p) {paglet\\\footnotesize image + passport};
  \draw[dashedarrow] (h3.south) -- node[event, below, sloped] {move} (p.east);
  \draw[dashedarrow] (p.west) -- node[event, below, sloped] {move} (h1.south);
\end{tikzpicture}
\caption{A mesh of three hosts. Every host holds a copy of the ledger and
  enforces the state it derives from it. Gossip spreads ledger records
  and host announcements; every host sends its own snapshot to every live
  host. An admin signs records on its own
  machine and brings them to any live host. An owner signs a passport for
  a paglet it launches; the passport moves with the paglet's image.}
\label{fig:mesh}
\end{figure}

\subsection{Hosts and the ledger}\label{sec:mesh-ledger}

Identities are Ed25519 keys in four roles: admins, hosts, owners and
signers (signers vouch for modules, Section~\ref{sec:movement}). A key's
ID is the hex form of its public key. Admin, owner and signer keys are
stored encrypted with a passphrase (Argon2id, XChaCha20-Poly1305); host
keys are stored unencrypted with owner-only file permissions, because a
host starts without a person.

The ledger is a grow-only set of signed records. A record is a canonical
MessagePack body with a set of signatures. Its ID is the SHA-256 of the
body, and signatures cover the ID, so a signature can be added later
without changing the record: this is how several admins co-sign a record
that needs a quorum. The first record is the genesis record: the mesh
name, the initial admins and the quorum, two numbers: how many admin
signatures an admin record needs, and how many a change of the admin set
needs (both 1 unless \code{mesh create} is given \code{-{}-quorum} and
\code{-{}-admin-quorum}; \code{admin-set} records change them). Its ID
is the mesh ID, with which a host opens its ledger copy as its trust
anchor.
Table~\ref{tab:ledger-records} lists the record types.

\begin{table}[htbp]
\centering\small
\begin{tabular}{@{}p{4.9cm}p{3.6cm}p{5.6cm}@{}}
\toprule
Record & Signed by & Content \tabularnewline
\midrule
\raggedright\code{genesis} & \raggedright every initial admin & \raggedright mesh name, admins, quorum \tabularnewline
\raggedright\code{admin-set} & \raggedright the admin-set quorum of its epoch & \raggedright admins added or removed, quorum \tabularnewline
\raggedright\code{host-enroll-request}, \code{owner-enroll-request} & \raggedright the key itself & \raggedright key, name; a host's labels \tabularnewline
\raggedright\code{host-enroll}, \code{owner-enroll} & \raggedright admin quorum & \raggedright key, name, labels (host) or groups (owner), the request answered \tabularnewline
\raggedright\code{host-remove}, \code{owner-remove} & \raggedright admin quorum & \raggedright key \tabularnewline
\raggedright\code{request-deny}, \code{revoke} & \raggedright admin quorum & \raggedright a request; a key, a record or a module \tabularnewline
\raggedright\code{module-sign} & \raggedright a signer & \raggedright module hash, name, version \tabularnewline
\raggedright\code{module-trust}, \code{module-policy} & \raggedright admin quorum & \raggedright modules or signers per trust class; whether roaming modules need trust \tabularnewline
\raggedright\code{root-residency} & \raggedright admin quorum & \raggedright where the content of a named root may go \tabularnewline
\raggedright\code{policy-rule} & \raggedright admin quorum & \raggedright decision, service, operations, match, scope \tabularnewline
\raggedright\code{grant-request} & \raggedright the paglet's host, or the owner & \raggedright principal, items, duration, reason \tabularnewline
\raggedright\code{grant} & \raggedright admin quorum, or a host under an allow rule & \raggedright principal, item, hosts, expiry \tabularnewline
\raggedright\code{grant-release}, \code{audit} & \raggedright a host & \raggedright a released grant; a decision \tabularnewline
\bottomrule
\end{tabular}
\caption{The ledger's record types (\code{cpp/host/include/paglets/mesh/ledger.hpp}).
  Records of unknown types are kept and replicated but have no effect, so
  newer record types pass through older hosts.}
\label{tab:ledger-records}
\end{table}

Adding a record checks only its form, its signatures and its mesh. What a
record means is decided when the state is derived from the whole set, in a
fixed order: by Lamport clock, then by ID. Hosts that hold the same
records derive the same state, whatever order the records arrived in.
Three rules close the gaps that a signed set leaves open. Admin changes
form a chain of epochs, and a signature counts only if its signer was an
admin of the epoch the record names. Because the signer chooses the
clock, the record that removes an admin lists the records that admin had
signed so far (its keep list); the removed admin's signatures count only
on those, so whatever it signs later, with any epoch and clock, has no
effect. Revocations win regardless of order: a revoked key is removed and
can never be enrolled again, and a revoked record has no effect. Records
without effect are listed with the reason (\code{LedgerState::ignored}),
and a digest of the derived state compares two hosts.

Replication is gossip over the host channels. A record that is new to a
host goes on at once to every host it knows except the sender; periodic
anti-entropy rounds exchange digests and inventories and repair what the
pushes missed. Every host holds the whole ledger and enforces the last
state it has, so allow and deny rules keep working while the mesh is
partitioned or no admin is online; when the partition heals, the records
merge. The ledger has no checkpoints yet: a host's copy grows with every
record.

A new host or owner signs an enrollment request with its own key, as proof
of possession, and an admin approves it from any host or enrolls the key
directly. An admin works on a local ledger copy and brings its records to
a live host with \code{remote push}; from there they spread. A host whose
key is not enrolled yet starts anyway and waits for its enrollment record.
Hosts find each other through a bootstrap contact (\code{--join URL}),
through signed host announcements that spread like records, and on a
local network through multicast beacons. A channel between two hosts
proves both host keys against the ledger, is encrypted end to end inside
HTTPS, and is refused between hosts of different mesh protocol versions or
paglet ABI major versions. Listing~\ref{lst:mesh-cli} builds a small
mesh: an admin creates it, enrolls a host and an owner, and adds a rule
under which the paglets of owners in the group \code{staff} must ask
before they read below \code{docs/} of the named root \code{data}; the
host binds that root to a directory when it starts
(\code{--root data=/srv/data}). The grant requests that such paglets later
make under the rule (Section~\ref{sec:mesh-policy}) are what
\code{remote requests} lists and \code{remote approve} decides;
\code{ledger audit} prints every decision the local copy holds.

\begin{lstlisting}[style=paglets-bash,float,caption={A mesh with one host and one owner, and a policy rule.},label={lst:mesh-cli}]
H=build/macos-arm64/host/paglets-host
$H keys init --role admin --name alice --out alice.key
$H keys init --role host --name lab-1 --out lab-1.key       # prints <lab-1-key-id>
$H keys init --role owner --name olga --out olga.key        # prints <olga-key-id>
$H mesh create --name lab --ledger ledger --admin alice.key
$H ledger enroll --ledger ledger --admin alice.key host <lab-1-key-id> lab-1 --label linux
$H ledger enroll --ledger ledger --admin alice.key owner <olga-key-id> olga --group staff
$H ledger rule --ledger ledger --admin alice.key --name "staff docs" --decision ask \
    --service files --op read --group staff --root data --path "docs/**"
$H serve --key lab-1.key --ledger ledger --state state-1 --listen 0.0.0.0:7443 \
    --advertise https://lab-1:7443 --root data=/srv/data
$H remote requests --connect https://lab-1:7443 --key alice.key --ledger ledger
$H remote approve  --connect https://lab-1:7443 --key alice.key --ledger ledger <request-id>
$H ledger audit --ledger ledger
\end{lstlisting}

\subsection{Passports}\label{sec:mesh-passports}

A paglet carries a passport. Its root is signed by the owner's key: the
mesh, the owner, the module hash, the paglet ID, a capability manifest,
and the issue and expiry times. \code{remote launch} issues one on the
owner's machine, valid for 24 hours unless \code{--hours} says otherwise,
and sends it to the host with the module and the start arguments. A child
or clone extends its creator's passport with a link signed by the host
that created it: the kind (\code{child} or \code{clone}), the parent and
the new paglet ID, the new module hash, the creating host, the issue
time, and the hash of the previous element, so a link cannot be moved to
another chain. A host verifies a
passport when it creates the paglet and when the paglet arrives from
another host: the signatures and the chain, the mesh, that the owner and
every linking host are enrolled and not revoked, that the passport has
not expired, and that its last element names the paglet and module in
question. The passport is stored beside the paglet and moves with its
image.

From the passport and the module the host forms the \emph{principal} that
every policy decision is made for: the paglet ID, the owner, the module
hash and the trust class (\code{roaming}, \code{resident} or
\code{system}). The host stamps the same fields into the sender record of
every delivered message, so a receiver can rely on them without a check
of its own. The manifest, a list of policy items the paglet declares in
advance, is evaluated by the same rules. The ledger accepts an
owner-signed \code{grant-request} for a paglet ID chosen in advance, so an
admin can approve the manifest before the paglet exists, and the grants
materialize when the paglet starts. No command of \code{paglets-host}
issues such a request yet, and \code{remote launch} issues passports with
an empty manifest.

\subsection{Policy rules and grants}\label{sec:mesh-policy}

Access to host resources is decided by rules in the ledger. The unit of
access is an \emph{item}: a service and its operations, and for
\code{files} a named root and a path below it, with rights (\code{read},
\code{write}, \code{create}, \code{delete}) in place of operations. A rule
names a service (or \code{*}) and operations, a decision (\emph{allow},
\emph{ask} or \emph{deny}), a match on the principal and the host, and for
files a scope. The match selects owners, owner groups, module hashes,
module signers, trust classes, and hosts by key or label; an absent member
matches everything. The scope lists roots and path patterns (\code{*},
\code{?}, \code{**}). A rule can bound the duration of the grants it gives
and carry a priority.

A host evaluates an item for a principal from its own ledger copy. Among
the rules that match the principal and the host, name the service and
cover all the operations and the scope, the most specific one wins:
highest priority, then most match members, then a scope over none, then
\emph{deny} over \emph{ask} over \emph{allow}, then ledger order. When no
rule matches, the decision is \emph{deny}. Every host evaluates the same
rules, so the decision does not depend on which host a paglet is on, only
on the host selectors of the rules.

The decision becomes records. Under an \emph{allow} rule the host signs a
\code{grant} that names the rule and itself as the only host; it expires
after the requested duration (1~s to 30 days), capped by the rule's
maximum. Other hosts accept such a grant only if the rule allows it for
that principal, so a host cannot mint access its rules do not give. Under
an \emph{ask} rule the host signs a \code{grant-request}; it is pending
until an admin, from any host, signs a \code{grant} that names it (with
\code{remote approve}, for every host or only the requesting one) or a
\code{request-deny}. A \emph{deny} leaves only the audit record. Every
decision writes an \code{audit} record signed by the host, and the admin's
approvals, denials and revocations are records themselves, so the audit
log is the ledger and is read from any host (\code{remote audit};
\code{ledger audit} prints the audit records a local copy holds). A grant
ends when it expires, when an admin revokes it, or when a host releases
it.

A paglet asks through the \code{grants} system paglet
(Listing~\ref{lst:grants-request}). The answer is \code{granted}, with the
capability as the first capability of the reply; \code{pending}, with the
request ID; or \code{denied}, with the deciding rule, if any. When an
approval arrives by gossip, the paglet's host \emph{materializes} the grant
and sends the capability in the message \code{grants.granted}
(Figure~\ref{fig:grant-flow}); a denial arrives as \code{grants.denied}. A
files item becomes a \code{dir} capability for the root and path with the
rights, the grant's ID and its expiry; an item of another service becomes
a non-transferable endpoint carrying the granted operations. The
capability records its grant, and the runtime refuses it with
\code{revoked} as soon as the record that revokes or releases the grant
arrives: every copy, derivative and transfer of the capability fails on
every host. An admin's grant belongs to the mesh, not to a host: it
materializes on every host its selector covers, also when the paglet
arrives there later, so it follows the paglet, while a host-local
capability does not (Section~\ref{sec:movement}). A grant a host derived
from an \emph{allow} rule names that host only; after a move the paglet
asks again, and the same rule allows again.

% snippet: cpp/common/include/paglets/services/grants.hpp#request [float,caption={The request and the answer of the \code{grants} service},label={lst:grants-request}]
\begin{lstlisting}[style=paglets-cpp,float,caption={The request and the answer of the \code{grants} service \src{cpp/common/include/paglets/services/grants.hpp}},label={lst:grants-request}]
struct Item {
    std::string service;           // files, server-info, ...
    std::vector<std::string> ops;  // files: rights (read, write, create, delete)
    std::string root;              // files: named root
    std::string path;              // files: directory below the root ("" for the root)
};

struct Request {
    Item item;
    std::int64_t duration_ms = 3'600'000;
    std::string reason;
};

enum class Status { granted, pending, denied };

struct Answer {
    Status status = Status::denied;
    Item item;
    std::string grant;    // hex ID of the grant record (granted)
    std::string request;  // hex ID of the request record (pending, and later answers)
    std::string rule;     // hex ID of the deciding rule, if any
    std::int64_t expires_ms = 0;
    std::string reason;
};
\end{lstlisting}
% /snippet

\begin{figure}[htbp]
\centering
\begin{tikzpicture}[x=1cm,y=1cm]
  \node[plain] (P) at (0,0) {paglet};
  \node[plain] (A) at (5.2,0) {host A: \code{grants}};
  \node[plain] (B) at (10.4,0) {host B: admin};
  \draw[lifeline] (P.south) -- (0,-7.0);
  \draw[lifeline] (A.south) -- (5.2,-7.0);
  \draw[lifeline] (B.south) -- (10.4,-7.0);
  \draw[msg] (0,-0.8) -- node[above] {\code{request}(item, duration, reason)} (5.2,-0.8);
  \node[event, anchor=west, align=left, text width=4.4cm] at (5.4,-1.6) {rule: \emph{ask}; records\\\code{grant-request}, \code{audit}};
  \draw[msg] (5.2,-2.6) -- node[above] {\code{Answer}: \code{pending}, request ID} (0,-2.6);
  \draw[msg] (5.2,-3.4) -- node[above] {gossip: \code{grant-request}} (10.4,-3.4);
  \node[event, anchor=east] at (10.2,-4.1) {\code{remote approve}: record \code{grant}};
  \draw[msg] (10.4,-4.9) -- node[above] {gossip: \code{grant}} (5.2,-4.9);
  \node[event, anchor=west] at (5.4,-5.6) {materialize a \code{dir} capability};
  \draw[msg] (5.2,-6.4) -- node[above] {\code{grants.granted} + capability} (0,-6.4);
\end{tikzpicture}
\caption{A grant request under an \emph{ask} rule. The admin decides on
  any host; the records travel by gossip; the paglet's host turns the
  approval into a capability.}
\label{fig:grant-flow}
\end{figure}

The pi example of Section~\ref{sec:mesh-compute} writes its digits below a
named root of its host. It asks for the directory with the patterns
helper \code{granted\_dir}, which sends the request and hands the
capability to a continuation (Listing~\ref{lst:pi-grant}). The helper
treats \code{pending} as a failure (\code{bad\_state}); a paglet that wants
to wait for an admin handles \code{grants.granted} itself.

% snippet: cpp/examples/pi/pi.cpp#grant [float,caption={The pi job asks for a directory of the named root it was started with},label={lst:pi-grant}]
\begin{lstlisting}[style=paglets-cpp,float,caption={The pi job asks for a directory of the named root it was started with \src{cpp/examples/pi/pi.cpp}},label={lst:pi-grant}]
pt::granted_dir(s.output_root, "", {"write", "create"}, "pi: digits", [this](paglets::Result<paglets::Capability> d) {
    if (!d) return fail(pt::error_text("files of " + params_.output_root, d.error()));
    dir_ = std::make_shared<paglets::Capability>(*d);
    const auto slash = params_.output_path.rfind('/');
    pt::make_dirs(*dir_, slash == std::string::npos ? std::string() : params_.output_path.substr(0, slash),
                  [this](paglets::Result<void> made) {
                      if (!made) return fail(pt::error_text("directory for " + params_.output_path, made.error()));
                      begin();
                  });
});
\end{lstlisting}
% /snippet

\subsection{System paglets}\label{sec:mesh-system-paglets}

A paglet runs in a sandbox: it cannot open files, create sockets or start
programs. Hosts lend them such resources through \emph{system paglets},
native services built into the host process. A system paglet is a paglet
for everybody else: it has the ID \code{system.<name>}, a mailbox, and
handles one message at a time; but it has no Wasm instance and no memory
image, and it never deactivates, moves or ends. Its handlers are trusted
host code without a time budget. Besides the ABI's imports
(Appendix~\ref{app:abi}: messages, capabilities, children and the
lifecycle, timers, the clock, random bytes and the log), these services
are the only way a paglet reaches anything on its host.
Table~\ref{tab:system-paglets} lists them.

\begin{table}[htbp]
\centering\small
\begin{tabular}{@{}l p{2.1cm} p{4.4cm} p{5.3cm}@{}}
\toprule
Service & Endpoint & Operations & Purpose \tabularnewline
\midrule
\code{files} & directory & \raggedright\code{list}, \code{stat}, \code{find}, \code{read}, \code{write}, \code{mkdir}, \code{move}, \code{remove}, \code{hash}, \code{share}, \code{save} & \raggedright files below a lent \code{dir} capability; \code{share} and \code{save} move files of any size between hosts as artifacts \tabularnewline
\code{grants} & default & \raggedright\code{request}, \code{status}, \code{release}, \code{list} & \raggedright access by the mesh policy \tabularnewline
\code{storage} & default & \raggedright\code{get}, \code{put}, \code{remove}, \code{list} & \raggedright the paglet's own key-value store, 16~MB \tabularnewline
\code{artifacts} & directory & \raggedright\code{put}, \code{get}, \code{stat} & \raggedright content-addressed blobs (SHA-256) that work on every host \tabularnewline
\code{pubsub} & directory & \raggedright\code{create}, \code{publish}, \code{subscribe}, \code{unsubscribe} & \raggedright topics on a host, reached by \code{topic} capabilities \tabularnewline
\code{server-info} & directory & \raggedright\code{summary}, \code{load}, \code{volumes}, \code{processes} & \raggedright the host's system, in one schema on every platform \tabularnewline
\code{directory} & default & \raggedright\code{lookup}, \code{publish}, \code{unpublish}, \code{list} & \raggedright services and paglets by name \tabularnewline
\code{user-info} & default & \raggedright\code{notify} & \raggedright notifications to the owner \tabularnewline
\code{locator} & default & \raggedright\code{locate}, \code{locate\_and\_pin}, \code{release} & \raggedright a paglet anywhere in the mesh; pins \tabularnewline
\code{mesh-info} & default & \raggedright\code{snapshot}, \code{landscape}, \code{select}, \code{find\_offers} & \raggedright every host's snapshot and offers \tabularnewline
\code{compute-slots} & default & \raggedright\code{request\_slot}, \code{release\_slot}, \code{status}, \code{candidates} & \raggedright admission of compute work \tabularnewline
\code{zombies} & default & \raggedright\code{join}, \code{peers}, \code{trace}, \code{ended}, \code{state}, \code{deliver}, \code{result}, \code{cancel}, \code{redirect}, \code{list}, \code{traces} & \raggedright the host's registry of zombie groups (Section~\ref{sec:zombies}) \tabularnewline
\code{web} & \raggedright directory, \code{--web} hosts & \raggedright\code{capabilities}, \code{fetch}, \code{extract\_text}, \code{download}, \code{search} & \raggedright mediated web access; downloads become artifacts \tabularnewline
\code{ai} & \raggedright directory, \code{--ai} hosts & \raggedright\code{capabilities}, \code{summarize}, \code{classify}, \code{extract}, \code{generate}, \code{embed}, \code{describe\_image} & \raggedright local inference through Ollama \tabularnewline
\bottomrule
\end{tabular}
\caption{The system paglets. \emph{Default}: every paglet is created with a
  non-transferable endpoint to the service. \emph{Directory}: a paglet looks
  the service up. For \code{server-info} the endpoint carries only the
  operations the policy allows; \code{files}, \code{artifacts} and
  \code{pubsub} act only on lent capabilities and need no policy;
  \code{web} and \code{ai} check the policy per request, against rules for
  the URL's host and path, or for the operation and the model.
  Every service also answers \code{describe} with its schema.}
\label{tab:system-paglets}
\end{table}

A paglet reaches a service in one of two ways. Services that act only on
the caller's own things (\code{directory}, \code{storage},
\code{user-info}, \code{grants}, \code{locator}, \code{mesh-info},
\code{compute-slots}, \code{zombies}) are installed as default endpoints
when the paglet is created (\code{paglets::service(name)} returns them).
Services that act only on capabilities the caller lends (\code{files},
\code{artifacts}, \code{pubsub}) are looked up in the \code{directory},
which returns an endpoint with every operation. Neither kind needs a
policy, because holding the capability is the permission.
\code{server-info} has ambient authority: it reads host state without any
capability. It is reached through the \code{directory}, and the endpoint
that a lookup returns carries only the operations the policy allows the
principal; the lookup fails with \code{denied} when it allows none. The
gateways \code{web} and \code{ai} are looked up the same way but decide
per request, against rules whose root and path are the URL's host and
path, or whose operation and scope name the task and the model. The
volume guard example works with \code{server-info}: a rule allows its
owner's paglets the operation \code{volumes}, and the guard looks the
service up, asks for the volumes and notifies its owner when one runs low
(Listing~\ref{lst:guard-check}). The first two calls of the listing
matter here: \code{lookup} returns the endpoint the policy allows, and
\code{volumes} succeeds only because the rule names that operation;
without the rule the lookup fails with \code{denied}. The rest is the
guard's bookkeeping of what it has reported.

% snippet: cpp/examples/guard/guard.cpp#check [float,caption={The volume guard reads the host's volumes through an endpoint whose operations the policy chose},label={lst:guard-check}]
\begin{lstlisting}[style=paglets-cpp,float,caption={The volume guard reads the host's volumes through an endpoint whose operations the policy chose \src{cpp/examples/guard/guard.cpp}},label={lst:guard-check}]
void check() {
    pt::lookup("server-info", [this](paglets::Result<paglets::Endpoint> ep) {
        if (!ep) return failed(pt::error_text("server-info", ep.error()));
        si::Client c{*ep};
        auto sent = c.volumes({}, [this](paglets::Result<si::Volumes> v, Message&) {
            if (!v) return failed(pt::error_text("volumes", v.error()));
            ++status_.checks;
            for (const auto& vol : v->volumes) {
                if (vol.total == 0) continue;
                const double free = 100.0 * static_cast<double>(vol.available) / static_cast<double>(vol.total);
                if (free >= watch_.min_free_percent) {
                    low_.erase(vol.mount_point);  // recovered
                    continue;
                }
                if (!low_.insert(vol.mount_point).second) continue;  // reported already
                status_.alerts.push_back(Alert{vol.mount_point, free, now_ms()});
                pt::notify(pt::Level::warning, "volume " + vol.mount_point + " is low on space",
                           std::to_string(static_cast<int>(free)) + "% free");
            }
            sleep();
        });
        if (!sent) failed(pt::error_text("volumes", sent.error()));
    });
}
\end{lstlisting}
% /snippet

Resource capabilities are the second half of the mechanism. A \code{dir}
capability names a root and a relative path and carries rights, and
optionally the grant it came from, an expiry and a use limit
(Section~\ref{sec:model-capabilities}). The guest never names a host
path: it lends the capability with a request to \code{files}, which acts
on what it is shown and never on its own authority; this is the
confused-deputy protection. A derived capability keeps the lineage and
the grant of its root, so one revocation reaches them all.
\code{files} resolves every path canonically
and refuses one that leads out of the root, also through symbolic links.
Content larger than a message travels as \emph{artifact} capabilities: a
few bytes that name the content by hash and the host that holds it. The
content moves from host to host only when a host needs it, within the
data residency of its root.

\subsection{Service contracts}\label{sec:mesh-contracts}

The system services are contracts of the form of
Section~\ref{sec:model-schemas}: a schema header in
\path{cpp/common/include/paglets/services/}, served on the host by
\code{services::ContractPaglet} and called from a guest through the
generated \code{Client}. Listings~\ref{lst:compute-types}
and~\ref{lst:compute-contract} show the request, the reply and the
contract of \code{compute-slots}.

% snippet: cpp/common/include/paglets/services/compute_slots.hpp#types [float,caption={The request and the reply of \code{request\_slot}},label={lst:compute-types}]
\begin{lstlisting}[style=paglets-cpp,float,caption={The request and the reply of \code{request\_slot} \src{cpp/common/include/paglets/services/compute\_slots.hpp}},label={lst:compute-types}]
struct SlotRequest {
    std::string job;  // the job it belongs to (for status)
    std::int64_t cores = 1;
    std::int64_t estimated_ms = 0;
};

enum class Verdict { run_now, queued, redirect, rejected };

struct Decision {
    Verdict verdict = Verdict::rejected;
    std::string lease;  // run_now
    std::int64_t cores = 0;
    std::string host;           // redirect: key ID of the host to go to
    std::string host_name;      // its name
    std::int64_t position = 0;  // queued: place in the queue (0: next)
    std::string reason;
};
\end{lstlisting}
% /snippet

% snippet: cpp/common/include/paglets/services/compute_slots.hpp#contract [float,caption={The contract of the \code{compute-slots} service},label={lst:compute-contract}]
\begin{lstlisting}[style=paglets-cpp,float,caption={The contract of the \code{compute-slots} service \src{cpp/common/include/paglets/services/compute\_slots.hpp}},label={lst:compute-contract}]
struct Contract {
    static constexpr std::string_view service = "compute-slots";
    Decision request_slot(const SlotRequest&);
    ReleaseReply release_slot(const ReleaseRequest&);
    Status status(const StatusRequest&);
    Candidates candidates(const CandidatesRequest&);
};
\end{lstlisting}
% /snippet

Host and guest embed the same JSON descriptor of the namespace,
byte-identical by construction, and every service answers \code{describe}
with it. (The forms of the tools in Section~\ref{sec:tools} come from a
module's interface description, a separate document that travels inside
the module.) A module names the services it uses in its build rule
(\code{paglets\_add\_module(... SERVICES compute\_slots mesh\_info)}) and
includes \code{<paglets/services/compute\_slots.gen.hpp>}; a caller
cannot tell a guest that serves a contract from a native system paglet.
Contracts need reflection on the host; a host built without it runs
paglets but has no services and no mesh.

\subsection{Mesh information and compute slots}\label{sec:mesh-compute}

Two services let a paglet see the mesh and get compute time on it.
\code{mesh-info} spreads a \emph{snapshot} of every host. Each host
samples itself every 2~s (CPUs, load average per CPU, CPU use, memory,
free space where it keeps state, paglets, compute slots, labels, offers,
and its links to the other hosts: up or down, since when, and why frames
fail) and sends the snapshot to every live host. A host accepts a snapshot
only from the host it describes, over the channel that proved its key; a
snapshot older than 20~s is not fresh, and after a minute it is dropped.
So every host answers for the whole mesh: \code{landscape} lists the hosts
with a fresh snapshot, \code{select} returns the hosts that fit a load,
memory, free-slot and label requirement, the best first, or, when none
fits, the least loaded hosts marked as a fallback, and
\code{find\_offers} the hosts whose services offer a feature, such as
\code{ai.summarize} with a model of a given context size. Offers are also
destinations: a transfer ticket \code{offer:ai.summarize}
(Table~\ref{tab:movement-tickets}) moves a paglet to such a host.

\code{compute-slots} is admission. A host has slots, by default one per
CPU (\code{serve --slots N} changes it); there is no central scheduler or
job queue, every host decides for itself and tells the others what it has
free. A job paglet asks the \code{compute-slots} of the host it is on and
gets a \code{Decision}: \code{run\_now} with a lease when slots are free
and nobody waits; \code{redirect} with a host that has free slots now and
no queue; \code{queued} with its position otherwise; \code{rejected} for
more cores than the host has, unless another host has them free now,
which is a \code{redirect}. The queue is served oldest-fit: the first
waiter, from the oldest, whose cores are free receives the message
\code{compute.granted}, which activates it if it deactivated while
waiting; a small waiter can pass a larger one ahead of it. Waiters behind
the first that have waited a second are offered to hosts that have free
slots and no queue (\code{compute.redirect}, at most four per round and
half the queue, with shadow reservations so that one host is not
flooded). The service never
moves a paglet: the paglet dispatches itself and asks again on arrival. A
lease ends when it is released, or when the paglet ends or leaves the
host.

Listing~\ref{lst:pi-worker} is a worker of the pi example, which computes
hexadecimal digits of pi in independent chunks, each in a clone of the job
on a host that \code{mesh-info} selected. The worker asks for one slot and
acts on the verdict; on a host without the service it runs at once. It
handles \code{compute.granted} by running and \code{compute.redirect} by
dispatching itself; \code{on\_arrived} asks again, and
\code{on\_move\_failed} asks the host it is still on. The job keeps the
state of every chunk and sends a chunk out again when it is not answered
in time, so the loss of a host costs only the chunks in flight there
(Section~\ref{sec:evaluation} reports a live run on three machines, and
the example's test, in which one of three hosts in one process is shut
down mid-run).

% snippet: cpp/examples/pi/pi.cpp#worker [float,caption={A pi worker asks its host for a compute slot},label={lst:pi-worker}]
\begin{lstlisting}[style=paglets-cpp,float,caption={A pi worker asks its host for a compute slot \src{cpp/examples/pi/pi.cpp}},label={lst:pi-worker}]
void ask_for_slot() {
    auto slots = paglets::service("compute-slots");
    if (!slots) return compute();  // a host without compute-slots: just run
    svc::compute_slots::Client client{*slots};
    svc::compute_slots::SlotRequest q;
    q.job = "pi";
    q.cores = 1;
    auto sent = client.request_slot(q, [this](paglets::Result<svc::compute_slots::Decision> d, Message&) {
        if (!d) return compute();
        switch (d->verdict) {
            case svc::compute_slots::Verdict::run_now: run(d->lease); break;
            case svc::compute_slots::Verdict::queued: break;  // compute.granted comes later
            case svc::compute_slots::Verdict::redirect:
                if (!paglets::dispatch(d->host)) compute();
                break;
            case svc::compute_slots::Verdict::rejected: compute(); break;
        }
    });
    if (!sent) compute();
}
\end{lstlisting}
% /snippet
